% GENERATED FILE. Edit canonical lessons, then run npm run generate:books.
% canonical-source-hash: fnv1a64:0b8e0c1c14399d3e
% canonical-lessons: RU-C62-hare, RU-C62-snake, RU-C62-hen, RU-C62-duck, RU-C62-goat

\chapter{Hare, Snake, Hen, Duck, Goat}
\label{ch:ru-hare-snake-hen-duck-goat}
\hlchaptermodality{62}

\begin{chapteropening}
\textbf{By the end of this chapter:} \emph{I can name a hare, a snake, a hen, a duck and a goat.}

Complete the last lesson of chapter 62: i can name a hare, a snake, a hen, a duck and a goat.
\end{chapteropening}

\section[záyats]{\ru{заяц} (záyats) — a hare}

\label{lesson:RU-C62-hare}



\begin{quote}
Before the new one: say the Russian for a bear, then the Russian for a fox.
\end{quote}

\subsection*{You'll want to know: \ru{заяц}}
\textbf{\ru{заяц}} (\emph{záyats}) — \textquotedblleft{}a hare\textquotedblright{}.

\begin{grammarlens}[title={What you've built}]
One more word for hare, snake, hen, duck, goat.
\end{grammarlens}

\subsection*{Your turn}
\begin{itemize}
\raggedright
  \item \emph{Say it:} \emph{záyats}
  \item \emph{Say it:} \emph{záyats}, once more
  \item \emph{Say it:} say \emph{záyats}
  \item \emph{Recall:} say the Russian for a bear, then the Russian for a fox, then say \emph{záyats} again
\end{itemize}

\begin{tcolorbox}[breakable,colback=teal!4,colframe=teal!35!black,title={Before you move on}]
What does \ru{заяц} mean? (\textquotedblleft{}A hare\textquotedblright{}.) Say it once more.
\end{tcolorbox}
\section[zmeyá]{\ru{змея} (zmeyá) — a snake}

\label{lesson:RU-C62-snake}



\begin{quote}
Before the new one: say the Russian for a fox, then the Russian for a hare.
\end{quote}

\subsection*{You'll want to know: \ru{змея}}
\textbf{\ru{змея}} (\emph{zmeyá}) — \textquotedblleft{}a snake\textquotedblright{}.

\begin{grammarlens}[title={What you've built}]
One more word for hare, snake, hen, duck, goat.
\end{grammarlens}

\subsection*{Your turn}
\begin{itemize}
\raggedright
  \item \emph{Say it:} \emph{zmeyá}
  \item \emph{Say it:} \emph{zmeyá}, once more
  \item \emph{Say it:} say \emph{zmeyá}
  \item \emph{Recall:} say the Russian for a fox, then the Russian for a hare, then say \emph{zmeyá} again
\end{itemize}

\begin{tcolorbox}[breakable,colback=teal!4,colframe=teal!35!black,title={Before you move on}]
What does \ru{змея} mean? (\textquotedblleft{}A snake\textquotedblright{}.) Say it once more.
\end{tcolorbox}
\section[kúritsa]{\ru{курица} (kúritsa) — a hen}

\label{lesson:RU-C62-hen}



\begin{quote}
Before the new one: say the Russian for a hare, then the Russian for a snake.
\end{quote}

\subsection*{You'll want to know: \ru{курица}}
\textbf{\ru{курица}} (\emph{kúritsa}) — \textquotedblleft{}a hen\textquotedblright{}. It lays the \textbf{\ru{яйцо}}.

\begin{grammarlens}[title={What you've built}]
One more word for hare, snake, hen, duck, goat.
\end{grammarlens}

\subsection*{Your turn}
\begin{itemize}
\raggedright
  \item \emph{Say it:} \emph{kúritsa}
  \item \emph{Say it:} \emph{kúritsa}, once more
  \item \emph{Say it:} say \emph{kúritsa}
  \item \emph{Recall:} say the Russian for a hare, then the Russian for a snake, then say \emph{kúritsa} again
\end{itemize}

\begin{tcolorbox}[breakable,colback=teal!4,colframe=teal!35!black,title={Before you move on}]
What does \ru{курица} mean? (\textquotedblleft{}A hen\textquotedblright{}.) Say it once more.
\end{tcolorbox}
\section[útka]{\ru{утка} (útka) — a duck}

\label{lesson:RU-C62-duck}



\begin{quote}
Before the new one: say the Russian for a snake, then the Russian for a hen.
\end{quote}

\subsection*{You'll want to know: \ru{утка}}
\textbf{\ru{утка}} (\emph{útka}) — \textquotedblleft{}a duck\textquotedblright{}.

\begin{grammarlens}[title={What you've built}]
One more word for hare, snake, hen, duck, goat.
\end{grammarlens}

\subsection*{Your turn}
\begin{itemize}
\raggedright
  \item \emph{Say it:} \emph{útka}
  \item \emph{Say it:} \emph{útka}, once more
  \item \emph{Say it:} say \emph{útka}
  \item \emph{Recall:} say the Russian for a snake, then the Russian for a hen, then say \emph{útka} again
\end{itemize}

\begin{tcolorbox}[breakable,colback=teal!4,colframe=teal!35!black,title={Before you move on}]
What does \ru{утка} mean? (\textquotedblleft{}A duck\textquotedblright{}.) Say it once more.
\end{tcolorbox}
\section[kozá]{\ru{коза} (kozá) — a goat}

\label{lesson:RU-C62-goat}



\begin{quote}
Before the new one: say the Russian for a hen, then the Russian for a duck.
\end{quote}

\subsection*{You'll want to know: \ru{коза}}
\textbf{\ru{коза}} (\emph{kozá}) — \textquotedblleft{}a goat\textquotedblright{}.

\begin{grammarlens}[title={What you've built}]
One more word for hare, snake, hen, duck, goat.
\end{grammarlens}

\subsection*{Your turn}
\begin{itemize}
\raggedright
  \item \emph{Say it:} \emph{kozá}
  \item \emph{Say it:} \emph{kozá}, once more
  \item \emph{Say it:} say \emph{kozá}
  \item \emph{Recall:} say the Russian for a hen, then the Russian for a duck, then say \emph{kozá} again
\end{itemize}

\begin{tcolorbox}[breakable,colback=teal!4,colframe=teal!35!black,title={Before you move on}]
What does \ru{коза} mean? (\textquotedblleft{}A goat\textquotedblright{}.) Say it once more.
\end{tcolorbox}
